\chapter{Як світ потрапляє всередину}
% full version: chapters/11_sensors.tex

Будь-який орган чуттів --- це перетворювач, як мікрофон чи фотоелемент. Світло, звук або молекула діють на особливий білок у клітині-датчику, білок відкриває «кран», тече струм, і наприкінці ланцюжка виходять клацання, які йдуть у мережу. Майже всі датчики викидають глутамат, тож входи машини підключаються просто до її телефонної мережі.

\section*{На межі фізики}

Датчики мозку працюють біля самої межі можливого. Датчик світла в темряві відповідає на одну-єдину частинку світла. Датчик звуку помічає зміщення, менше за нанометр, --- менше за розмір кількох атомів. У напівтемряві точність обмежує вже не око, а саме світло: частинки приходять випадково, і щоб розрізнити відтінок, їх треба накопичити більше. Тому в сутінках ми бачимо повільніше --- око довше накопичує світло, як фотоапарат із довгою витримкою.

\section*{Автоматична експозиція}

Освітленість від зоряної ночі до сонячного снігу змінюється в мільярди разів, гучність звуку --- ще сильніше. А частота клацань змінюється від нуля до кількох сотень на секунду. Втиснути одне в інше напряму неможливо. Розв'язок той самий, що у фотоапарата: автоматична експозиція. Датчик весь час підлаштовує чутливість під середній рівень навколо й повідомляє не яскравість, а яскравість \emph{порівняно з тлом}. Стале з часом перестає викликати відповідь, а кожна зміна --- викликає. Входи машини від самого початку говорять про зміни.

\pencilfig{125}{%
% light rises in two tenfold steps, then drops a hundredfold; sensor rate = baseline + gain * (log light - adapted level),
% the adapted level follows log light with a 0.7 s time constant; rate clipped at zero
\draw[penlite=pcAmber] (0,1.75) -- (1.4,1.75) -- (1.4,2.1) -- (4.2,2.1) -- (4.2,2.45) -- (6.3,2.45) -- (6.3,1.75) -- (8.4,1.75);
\node[lbl, text=pcAmber!70!black, anchor=east] at (-0.05,2.0) {світло};
\node[lbl, anchor=south] at (2.8,2.12) {яскравіше в 10}; \node[lbl, anchor=south] at (5.25,2.47) {ще в 10};
\node[lbl, anchor=south] at (7.35,1.77) {знову темно};
\draw[pen] (0,0) -- (8.4,0);
\draw[pcGray, densely dashed, line width=0.5pt] (0,0.32) -- (8.4,0.32);
\draw[penlite=pcBlue] plot coordinates {(0.00,0.32) (1.26,0.32) (1.40,1.02) (1.54,0.85) (1.68,0.72) (1.82,0.62) (1.96,0.54) (2.10,0.49) (2.24,0.45) (2.38,0.41) (2.52,0.39) (2.66,0.37) (2.80,0.36) (3.08,0.34) (3.50,0.33) (4.06,0.32) (4.20,1.03) (4.34,0.85) (4.48,0.72) (4.62,0.62) (4.76,0.54) (4.90,0.49) (5.04,0.45) (5.18,0.41) (5.32,0.39) (5.46,0.37) (5.60,0.36) (5.88,0.34) (6.16,0.33) (6.30,0.00) (7.00,0.00) (7.14,0.07) (7.28,0.13) (7.42,0.18) (7.56,0.22) (7.70,0.24) (7.84,0.26) (7.98,0.28) (8.12,0.29) (8.40,0.30)};
\node[lbl, text=pcBlue, anchor=east] at (-0.05,0.32) {відповідь};
\node[lbl, anchor=north] at (6.65,-0.02) {мовчить};
}{Світло стає яскравішим сходинками, кожна вдесятеро. Датчик відповідає на кожну зміну сплеском і швидко повертається до попереднього рівня. Коли знову темніє, він на час замовкає.}

Інженери повторили цей прийом у подієвих камерах. Така камера не знімає кадри за годинником: кожен її піксель сам повідомляє «стало світліше» чи «стало темніше», коли щось змінилося. Це та сама пара проводів «так» і «ні» з другого розділу, тільки в кремнії.

\section*{Око стискає картинку}

В оці близько ста мільйонів датчиків світла, а проводів, що везуть зображення в мозок, --- близько мільйона. Перш ніж картинка покине око, її стискають приблизно в сто разів. Рівні ділянки майже нічого не коштують, передаються краї та зміни. За оцінкою, зробленою на очах тварин і перерахованою на людину, око передає в мозок порядку десяти мільйонів бітів на секунду --- як старий мережевий кабель.

\section*{Вухо --- це піаніно}

Вухо розкладає звук за частотами, як інженер поставив би набір фільтрів. Пружна перетинка у вусі в різних місцях відгукується на різні частоти, і датчики на кожній ділянці слухають свою «клавішу». Номер датчика, що спрацював, --- це висота тону.

\pencilfig{11}{\pencilart{11}}{Вухо влаштоване як піаніно навпаки: кожна клавіша слухає свою ноту.}

Ба більше, вухо не пасивне. Частина його клітин сама розгойдує перетинку в такт звуку й підсилює тихі звуки в тисячі разів, а гучні --- майже не підсилює. Це підсилювач біля самої межі самозбудження: там він особливо чутливий. А на низьких частотах нейрони ще й клацають у такт хвилі, зберігаючи точний час звуку, --- за ним мозок і знаходить напрямок на джерело. Що вища частота, то гірше нейрони за нею встигають, і лишається тільки код місцем.

\section*{Інші чуття}

Майже в кожному чутті впізнається прийом із попередніх розділів. Температуру передають два проводи, окремі для тепла й для холоду. Дотик користується датчиками, які відповідають на доторк і одразу замовкають, і датчиками, які тримають відповідь, поки є тиск. А нюх влаштований як штрихкод: типів нюхових датчиків у людини близько чотирьохсот, і запах --- це набір типів, що спрацювали. Кількість таких наборів астрономічна.

\fullversion{11}
